\section{Arquitetura proposta}

O sistema segue a arquitetura \textbf{mestre/trabalhador}
(\emph{master-worker})~\cite{masterworker}, no modelo de paralelismo de
dados popularizado pelo MapReduce~\cite{mapreduce}
(Figura~\ref{fig:arquitetura}):

\begin{enumerate}
  \item o \textbf{mestre} divide os dados em pedaços (\emph{chunks}) e
    distribui as tarefas;
  \item os \textbf{workers} processam seus pedaços em paralelo e devolvem
    resultados parciais pequenos (uma soma, uma contagem), nunca os dados
    brutos;
  \item o mestre \textbf{agrega} os resultados parciais e produz a resposta
    final.
\end{enumerate}

\begin{figure}[H]
  \centering
  \begin{tikzpicture}[
      no/.style={draw, rounded corners, minimum width=2.6cm,
                 minimum height=0.9cm, align=center},
      seta/.style={-{Stealth[length=2.5mm]}, thick},
      node distance=1.1cm and 0.6cm]
    \node[no, fill=blue!10] (dados) {Dados de entrada};
    \node[no, fill=orange!20, below=of dados] (mestre)
      {\textbf{Mestre}\\\small divide os dados};
    \node[no, fill=green!15, below=of mestre] (w2) {Worker 2};
    \node[no, fill=green!15, left=of w2] (w1) {Worker 1};
    \node[no, fill=green!15, right=of w2] (w3) {Worker N};
    \node[no, fill=orange!20, below=of w2] (agrega)
      {\textbf{Mestre}\\\small agrega os parciais};
    \node[no, fill=blue!10, below=of agrega] (resp) {Resposta final};

    \draw[seta] (dados) -- (mestre);
    \foreach \w in {w1, w2, w3} {
      \draw[seta] (mestre) -- node[pos=0.5, font=\scriptsize, fill=white]
        {pedaço} (\w);
      \draw[seta] (\w) -- (agrega);
    }
    \draw[seta] (agrega) -- (resp);
  \end{tikzpicture}
  \caption{Fluxo da arquitetura mestre/workers.}
  \label{fig:arquitetura}
\end{figure}

Há duas formas de entregar os pedaços aos workers, e o grupo vai comparar
as duas:

\begin{description}
  \item[Estática:] os dados são divididos em $N$ partes iguais, uma por
    worker. É simples, mas se um worker for mais lento, todos esperam por
    ele.
  \item[Dinâmica (\emph{self-scheduling}):] os dados são divididos em
    muitos pedaços pequenos, e o mestre envia o próximo pedaço a cada
    worker que termina. Workers mais rápidos processam mais pedaços, e
    ninguém fica ocioso~\cite{calvinmpi}.
\end{description}

Diferentemente do cenário de balanceamento de carga, que distribui
\emph{requisições} contínuas com foco em tempo de resposta, este cenário
distribui \emph{pedaços de dados} de um trabalho com começo e fim
(processamento em lote), com foco em terminar mais rápido.
